%Data institusi

%\input{Dozentdaten}

%\renewcommand{\fachbereich}{Bidang studi}

%\renewcommand{\dozent}{Prof. Dr. . }

%Data ujian

\renewcommand{\klausurgebiet}{ }

\renewcommand{\klausurtyp}{ }

\renewcommand{\klausurdatum}{. 20}

\klausurvorspann {\fachbereich} {\klausurdatum} {\dozent} {\klausurgebiet} {\klausurtyp}

%Data untuk tabel poin berikut

\renewcommand{\aeins}{  3  }

\renewcommand{\azwei}{  3   }

\renewcommand{\adrei}{  5  }

\renewcommand{\avier}{  2  }

\renewcommand{\afuenf}{  4  }

\renewcommand{\asechs}{  8  }

\renewcommand{\asieben}{  4  }

\renewcommand{\aacht}{  5  }

\renewcommand{\aneun}{  3  }

\renewcommand{\azehn}{  2  }

\renewcommand{\aelf}{  6  }

\renewcommand{\azwoelf}{  7  }

\renewcommand{\adreizehn}{  4   }

\renewcommand{\avierzehn}{  8  }

\renewcommand{\afuenfzehn}{  64  }

\renewcommand{\asechzehn}{    }

\renewcommand{\asiebzehn}{    }

\renewcommand{\aachtzehn}{    }

\renewcommand{\aneunzehn}{    }

\renewcommand{\azwanzig}{    }

\renewcommand{\aeinundzwanzig}{    }

\renewcommand{\azweiundzwanzig}{    }

\renewcommand{\adreiundzwanzig}{    }

\renewcommand{\avierundzwanzig}{    }

\renewcommand{\afuenfundzwanzig}{    }

\renewcommand{\asechsundzwanzig}{    }

\punktetabellevierzehn

\klausurnote

\newpage

\setcounter{section}{0}

__NOEDITSECTION__

\inputaufgabeklausurloesung
{3}

{

Definisikan istilah-istilah berikut
\zusatzklammer {yang dicetak miring} {} {}.
\aufzaehlungsechs{\stichwort {turunan tangensial kedua} {}
\zusatzklammer {atau
\stichwort {percepatan tangensial} {}} {} {}
\zusatzklammer {sebagai pemetaan ke $\R^n$} {} {}
dari suatu
\definitionsverweis {kurva yang dua kali terdiferensialkan secara kontinu}{}{}
\maabbdisp {\gamma} {I} {Y
} {}
yang bernilai dalam suatu
\definitionsverweis {hipermuka terdiferensialkan}{}{}

\mavergleichskette
{\vergleichskette
{ Y
}
{ \subseteq }{ \R^n
}
{  }{
}
{    }{
}
{   }{
}
}
{}{}{.}

}{Suatu ruang topologis $X$ yang
\stichwort {terhubung} {}.

}{\stichwort {Pemetaan singgung di titik} {}

\mavergleichskette
{\vergleichskette
{  P
}
{ \in }{ M
}
{  }{}
{    }{}
{   }{}
}
{}{}{}
untuk suatu pemetaan diferensiabel
\maabbdisp {\varphi} { M } { N
} {,}
dengan
\mathkor {} {M} {dan} {N} {}
manifold diferensiabel.

}{\stichwort {Produk} {}
dari manifold diferensiabel
\mathkor {} {M} {dan} {N} {.}

}{\stichwort {Ukuran volume} {}
yang terkait dengan bentuk volume positif $\omega$ pada suatu manifold diferensiabel berdimensi $n$.

}{Suatu
\stichwort {kurva geodesik} {}
pada manifold Riemann $M$.
}

}
{

\aufzaehlungsechs{Pemetaan
\maabbeledisp {} {I} { TY
} {t} { \pi_P( \gamma^{\prime \prime} (t))
} {,}
dengan $\pi_P$ menyatakan proyeksi ortogonal
\maabbdisp {} {\R^n } { T_PY
} {}
disebut
turunan tangensial kedua
\zusatzklammer {atau
\definitionswort {percepatan tangensial}{}} {} {}
dari $\gamma$.
}{Suatu
\definitionsverweis {ruang topologis}{}{}
$X$ disebut terhubung jika tepat ada dua himpunan bagian di $X$
\zusatzklammer {yaitu $\emptyset$ dan seluruh ruang

\mavergleichskettek
{\vergleichskettek
{  X
}
{ \neq }{ \emptyset
}
{  }{
}
{    }{
}
{   }{
}
}
{}{}{}} {} {,}
yang sekaligus
\definitionsverweis {terbuka}{}{}
dan
\definitionsverweis {tertutup}{}{}.
}{Yang dimaksud dengan pemetaan singgung di titik $P$ adalah pemetaan
\maabbeledisp {} {T_PM} {T_{\varphi(P)}N
} {[\gamma]} {[\varphi \circ \gamma]
} {.}
}{Misalkan
\mathkor {} {(U_i,U_i',\alpha_i, i \in I)} {dan} {(V_j,V_j',\beta_j, j \in J)} {}
adalah
\definitionsverweis {atlas}{}{}
dari
\mathkor {} {M} {dan} {N} {.}
Maka
\definitionsverweis {ruang produk}{}{}

\mathl{M \times N}{} yang dilengkapi dengan
\definitionsverweis {peta}{}{}
\maabbdisp {\alpha_i \times \beta_j} {U_i \times V_j } { U_i' \times V_j'
} {}
\zusatzklammer {dengan \mathlk{(i,j) \in I \times J}{} dan \mathlk{U_i' \times V_j' \subseteq \R^m \times \R^n}{}} {} {}
disebut produk dari manifold
\mathkor {} {M} {dan} {N} {.}
}{Untuk suatu
\definitionsverweis {himpunan Borel}{}{}

\mavergleichskette
{\vergleichskette
{ T
}
{ \subseteq }{ M
}
{  }{
}
{    }{
}
{   }{
}
}
{}{}{}
ukuran $T$ terhadap $\omega$ didefinisikan melalui suatu
\definitionsverweis {dekomposisi}{}{}
\definitionsverweis {terhitung}{}{}

\mavergleichskette
{\vergleichskette
{ T
}
{ = }{  \biguplus_{i \in I} T_i
}
{  }{
}
{    }{
}
{   }{
}
}
{}{}{}
\zusatzklammer {dengan

\mavergleichskettek
{\vergleichskettek
{ T_i
}
{ \subseteq }{U_i
}
{  }{
}
{    }{
}
{   }{
}
}
{}{}{}
suatu domain peta terbuka dan

\mavergleichskettek
{\vergleichskettek
{  \alpha_{i*} \omega
}
{ = }{ f_i dx_1 \wedge \ldots \wedge dx_n
}
{  }{
}
{    }{
}
{   }{
}
}
{}{}{}
berlaku} {} {}

\mavergleichskettedisp
{\vergleichskette
{
\int_{ T  }  \omega
}
{ =} { \sum_{i \in I} \int_{ \alpha_i(T_i)  }  f_i   \,  d \lambda^n
}
{ } {
}
{ } {
}
{ } {
}
}
{}{}{}
sebagaimana di atas.
}{Suatu
\definitionsverweis {kurva diferensiabel}{}{}
\maabbdisp {\gamma} {I} {M
} {}
disebut
kurva geodesik
jika

\mavergleichskettedisp
{\vergleichskette
{ \nabla_{d/dt}  \gamma'
}
{ =} { 0
}
{ } {
}
{ } {
}
{ } {
}
}
{}{}{}
berlaku pada $I$.
}

 }

__NOEDITSECTION__

\inputaufgabeklausurloesung
{3}

{

Nyatakan teorema-teorema berikut.
\aufzaehlungdrei{\stichwort {Teorema tentang kelengkungan kurva bidang yang diberikan secara implisit} {}

\mavergleichskette
{\vergleichskette
{  Y
}
{ \subseteq }{  \R^2
}
{  }{
}
{    }{
}
{   }{
}
}
{}{}{.}}{\stichwort {Rumus transformasi untuk bentuk volume positif pada manifold} {.}}{Teorema tentang kelengkungan Gauss dan kelengkungan seksional pada suatu permukaan diferensiabel

\mavergleichskette
{\vergleichskette
{  Y
}
{ \subseteq }{  \R^3
}
{  }{
}
{    }{
}
{   }{
}
}
{}{}{}}

}
{

\aufzaehlungdrei{\faktsituation {Misalkan

\mavergleichskette
{\vergleichskette
{ W
}
{ \subseteq }{ \R^2
}
{  }{
}
{    }{
}
{   }{
}
}
{}{}{}
\definitionsverweis {terbuka}{}{,}
\maabb {h} { W } { \R
} {}
adalah suatu
\definitionsverweis {fungsi yang dua kali terdiferensialkan secara kontinu}{}{}
dan

\mavergleichskette
{\vergleichskette
{ Y
}
{ = }{ h^{-1}(c)
}
{  }{
}
{    }{
}
{   }{
}
}
{}{}{}
adalah
\definitionsverweis {serat}{}{}
di atas

\mavergleichskette
{\vergleichskette
{ c
}
{ \in }{ \R
}
{  }{
}
{    }{
}
{   }{
}
}
{}{}{,}
dengan $h$
\definitionsverweis {reguler}{}{}
di setiap titik $Y$. Misalkan
\maabbdisp {N} { U } {\R^2
} {}
yang pada suatu lingkungan terbuka

\mavergleichskette
{\vergleichskette
{ Y
}
{ \subseteq }{ U
}
{  }{
}
{    }{
}
{   }{
}
}
{}{}{}
didefinisikan sebagai
\definitionsverweis {medan normal satuan}{}{}
untuk $Y$, dan misalkan

\mavergleichskette
{\vergleichskette
{ P
}
{ \in }{ Y
}
{  }{
}
{    }{
}
{   }{
}
}
{}{}{.}}
\faktfolgerung {Maka, untuk setiap vektor singgung

\mavergleichskette
{\vergleichskette
{ v
}
{ \in }{ T_PY
}
{  }{
}
{    }{
}
{   }{
}
}
{}{}{}

\mavergleichskettedisp
{\vergleichskette
{ \left(  D N   \right)_{ P } \left(  v  \right)
}
{ =} { - \kappa(P) v
}
{ } {
}
{ } {
}
{ } {
}
}
{}{}{,}
dengan $\kappa(P)$ adalah
\definitionsverweis {kelengkungan}{}{}
dari suatu parametrisasi panjang busur $Y$ yang sesuai dengan orientasi yang ditentukan oleh $v, N(P)$.}
\faktzusatz {}
\faktzusatz {}}{\faktsituation {Misalkan $M$ adalah suatu
\definitionsverweis {manifold diferensiabel}{}{}
dengan
\definitionsverweis {basis topologi terhitung}{}{}
dan misalkan $\omega$ adalah suatu
\definitionsverweis {bentuk volume positif}{}{}
pada $M$. Misalkan
\maabbdisp {\varphi} { L } { M
} {}
adalah suatu
\definitionsverweis {difeomorfisme}{}{}
dari manifold $L$ dan

\mavergleichskette
{\vergleichskette
{ S
}
{ \subseteq }{ L
}
{  }{
}
{    }{
}
{   }{
}
}
{}{}{}
adalah suatu
\definitionsverweis {himpunan bagian terukur}{}{.}}
\faktfolgerung {Maka

\mavergleichskettedisp
{\vergleichskette
{ \int_S \varphi^* \omega
}
{ =} { \int_{ \varphi(S)} \omega
}
{ } {
}
{ } {
}
{ } {
}
}
{}{}{.}}
\faktzusatz {}
\faktzusatz {}}{\faktsituation {Misalkan

\mavergleichskette
{\vergleichskette
{ W
}
{ \subseteq }{ \R^3
}
{  }{
}
{    }{
}
{   }{
}
}
{}{}{}
\definitionsverweis {terbuka}{}{}
dan misalkan
\maabb {h} { W } { \R
} {}
adalah suatu
\definitionsverweis {fungsi yang dua kali terdiferensialkan secara kontinu}{}{}
dan

\mavergleichskette
{\vergleichskette
{ Y
}
{ = }{ h^{-1}(c)
}
{  }{
}
{    }{
}
{   }{
}
}
{}{}{}
adalah serat di atas

\mavergleichskette
{\vergleichskette
{ c
}
{ \in }{ \R
}
{  }{
}
{    }{
}
{   }{
}
}
{}{}{,}
dengan $h$
\definitionsverweis {reguler}{}{}
di setiap titik $Y$.}
\faktfolgerung {Maka
\definitionsverweis {kelengkungan Gauss}{}{}
dari $Y$ sama dengan
\definitionsverweis {kelengkungan seksional}{}{}
dari $Y$.}
\faktzusatz {}
\faktzusatz {}}

 }

__NOEDITSECTION__

\inputaufgabeklausurloesung
{5}

{

Misalkan

\mavergleichskette
{\vergleichskette
{ G
}
{ \subseteq }{ \R^n
}
{  }{
}
{    }{
}
{   }{
}
}
{}{}{}
\definitionsverweis {terbuka}{}{}
dan
\maabbdisp {\varphi} { G } { \R^m
} {}
adalah
\definitionsverweis {pemetaan yang terdiferensialkan secara kontinu}{}{,}
yang pada titik

\mavergleichskette
{\vergleichskette
{ P
}
{ \in }{ G
}
{  }{
}
{    }{
}
{   }{
}
}
{}{}{}
memiliki
\definitionsverweis {diferensial total}{}{}
yang surjektif. Misalkan

\mavergleichskette
{\vergleichskette
{ v
}
{ \in }{ T_PZ
}
{  }{
}
{    }{
}
{   }{
}
}
{}{}{}
adalah suatu vektor dalam
\definitionsverweis {ruang singgung serat}{}{}
$Z$ dari $\varphi$ yang melalui $P$. Tunjukkan bahwa terdapat
\definitionsverweis {kurva yang terdiferensialkan secara kontinu}{}{}
\maabbdisp {\gamma} { ]- \epsilon, \epsilon[ } { Z
} {}
\zusatzklammer {untuk suatu

\mavergleichskettek
{\vergleichskettek
{ \epsilon
}
{ > }{ 0
}
{  }{
}
{    }{
}
{   }{
}
}
{}{}{}
yang sesuai} {} {}
dengan

\mavergleichskette
{\vergleichskette
{ \gamma(0)
}
{ = }{ P
}
{  }{
}
{    }{
}
{   }{
}
}
{}{}{}
dan

\mavergleichskettedisp
{\vergleichskette
{ \gamma'(0)
}
{ =} { v
}
{ } {
}
{ } {
}
{ } {
}
}
{}{}{}.

}
{

Berlaku

\mavergleichskette
{\vergleichskette
{v
}
{ \in }{ T_PZ
}
{ = }{ \operatorname{kern}  \left(D\varphi\right)_{P}
}
{    }{
}
{   }{
}
}
{}{}{.}
Berdasarkan
[[Satz über implizite Abbildungen/R/Fakt|Satz über implizite Abbildungen]],
terdapat suatu himpunan terbuka

\mathbed {P \in W} {}
 {W \subseteq G} {}
 {} {} {} {,}
suatu himpunan terbuka

\mavergleichskette
{\vergleichskette
{ V
}
{ \subseteq }{ \R^{n-m}
}
{  }{
}
{    }{
}
{   }{
}
}
{}{}{}
dan suatu pemetaan yang terdiferensialkan secara kontinu
\maabbdisp {\psi} { V } { W
} {}
sedemikian sehingga

\mavergleichskette
{\vergleichskette
{ \psi(V)
}
{ \subseteq }{ Z \cap W
}
{  }{
}
{    }{
}
{   }{
}
}
{}{}{}
berlaku dan $\psi$ menginduksi suatu
\definitionsverweis {bijeksi}{}{}
\maabbdisp {\psi} { V } { Z \cap W
} {}
tersebut. Misalkan

\mavergleichskette
{\vergleichskette
{ Q
}
{ \in }{ V
}
{  }{
}
{    }{
}
{   }{
}
}
{}{}{}
adalah titik yang memenuhi

\mavergleichskette
{\vergleichskette
{ \psi(Q)
}
{ = }{ P
}
{  }{
}
{    }{
}
{   }{
}
}
{}{}{.}
Pemetaan $\psi$ di setiap titik

\mavergleichskette
{\vergleichskette
{ Q
}
{ \in }{ V
}
{  }{
}
{    }{
}
{   }{
}
}
{}{}{}
bersifat
\definitionsverweis {reguler}{}{}
dan untuk
\definitionsverweis {diferensial total}{}{}
dari $\psi$ berlaku

\mavergleichskettedisp
{\vergleichskette
{ \left(D \varphi \right)_{ P }  \circ \left(D \psi \right)_{ Q }
}
{ =} { 0
}
{ } {
}
{ } {
}
{ } {
}
}
{}{}{,}
sehingga

\mavergleichskettedisp
{\vergleichskette
{ \operatorname{bild}  \left(D \psi\right)_{ Q }
}
{ =} { \operatorname{kern}  \left(D\varphi \right)_{ P }
}
{ =} { T_P Z
}
{ } {
}
{ } {
}
}
{}{}{.}
Karena regularitas $\psi$ di $Q$, pemetaan
\maabbdisp {\left(D \psi \right)_{ Q }} { \R^{n-m}  } { \R^n
} {}
bersifat injektif dan pemetaan
\maabbdisp {\left(D \psi \right)_{ Q }} { \R^{n-m}  } { T_PZ
} {}
bersifat bijektif. Misalkan

\mavergleichskette
{\vergleichskette
{ u
}
{ \in }{ \R^{n-m}
}
{  }{
}
{    }{
}
{   }{
}
}
{}{}{}
adalah prapeta dari $v$ dan misalkan
\maabbeledisp {\theta} { ]- \epsilon, \epsilon[ } { \R^{n-m}
} { t } { Q+tu
} {,}
dengan $\epsilon$ dipilih cukup kecil sehingga citranya seluruhnya termuat dalam $V$. Maka
\maabbdisp {\gamma = \psi \circ \theta} { ]- \epsilon, \epsilon[  } { \R^n
} {}
memenuhi

\mavergleichskettedisp
{\vergleichskette
{ \gamma(0)
}
{ =} { \psi( \theta (0))
}
{ =} { \psi(Q)
}
{ =} { P
}
{ } {
}
}
{}{}{}
dan

\mavergleichskettedisp
{\vergleichskette
{ \gamma' (0)
}
{ =} { \left(D \psi\right)_{Q} ( \theta' (0))
}
{ =} { \left(D \psi\right)_{Q} ( u)
}
{ =} { v
}
{ } {
}
}
{}{}{.}

 }

__NOEDITSECTION__

\inputaufgabeklausurloesung
{2}

{

Misalkan
\mathkor {} {r} {dan} {R} {}
adalah bilangan real positif dengan

\mavergleichskette
{\vergleichskette
{ 0
}
{ < }{ r
}
{ < }{ R
}
{    }{
}
{   }{
}
}
{}{}{.}
Tentukan suatu
\definitionsverweis {medan normal satuan}{}{}
untuk torus
\zusatzklammer {tertanam} {} {}

\mavergleichskettedisp
{\vergleichskette
{ T
}
{ =} { \left\{ (x,y,z) \in \R^3  \mid   \left(  \sqrt{x^2+y^2} -R  \right)^2 +z^2  = r^2         \right\}
}
{ } {
}
{ } {
}
{ } {
}
}
{}{}{.}

}
{

Berlaku

\mavergleichskettedisp
{\vergleichskette
{ h(x,y,z)
}
{ =} { x^2+y^2 -2 \sqrt{x^2+y^2}R+z^2 + R^2
}
{ } {
}
{ } {
}
{ } {
}
}
{}{}{}
dan

\mavergleichskette
{\vergleichskette
{ T
}
{ = }{ h^{-1} (r^2)
}
{  }{
}
{    }{
}
{   }{
}
}
{}{}{.}
Gradien dari $h$ adalah

\mavergleichskettedisp
{\vergleichskette
{
\operatorname{Grad} \,  h
}
{ =} { 2 \begin{pmatrix}   x   -  xR \left(  x^2+y^2  \right)^{- { \frac{   1  }{  2 } } }  \\ y   -  y R \left(  x^2+y^2  \right)^{- { \frac{   1  }{  2 } } } \\ z   \end{pmatrix}
}
{ } {
}
{ } {
}
{ } {
}
}
{}{}{.}
Dengan demikian,
\definitionsverweis {medan normal satuan}{}{}
adalah

\mavergleichskettedisp
{\vergleichskette
{ N(x,y,z)
}
{ =} { { \frac{   1  }{  \sqrt{ \left(   x   -  xR \left(  x^2+y^2  \right)^{- { \frac{   1  }{  2 } } }  \right)^2 + \left(  y   -  y R \left(  x^2+y^2  \right)^{- { \frac{   1  }{  2 } } }  \right)^2   +z^2   }  } } \begin{pmatrix}   x   -  xR \left(  x^2+y^2  \right)^{- { \frac{   1  }{  2 } } }  \\ y   -  y R \left(  x^2+y^2  \right)^{- { \frac{   1  }{  2 } } } \\ z   \end{pmatrix}
}
{ } {
}
{ } {
}
{ } {
}
}
{}{}{.}

 }

__NOEDITSECTION__

\inputaufgabeklausurloesung
{4}

{

Jelaskan peran teorema tentang pemetaan implisit dalam pengembangan teori manifold.

}
{  }

__NOEDITSECTION__

\inputaufgabeklausurloesung
{8}

{

Buktikan teorema tentang sifat adjoin-diri (self-adjoint) dari pemetaan Weingarten.

}
{

Untuk vektor-vektor

\mavergleichskette
{\vergleichskette
{ v,w
}
{ \in }{ T_PY
}
{  }{
}
{    }{
}
{   }{
}
}
{}{}{}
harus ditunjukkan bahwa

\mavergleichskettedisp
{\vergleichskette
{ \left\langle  v  ,  L_P(w)  \right\rangle
}
{ =} { \left\langle L_P(v) ,  w  \right\rangle
}
{ } {
}
{ } {
}
{ } {
}
}
{}{}{}
berlaku. Dengan

\mavergleichskette
{\vergleichskette
{ N(P)
}
{ = }{ { \frac{
\operatorname{Grad} \,  h  (  P )  }{  \Vert {
\operatorname{Grad} \,  h  (  P ) } \Vert  } }
}
{  }{
}
{    }{
}
{   }{
}
}
{}{}{}
berdasarkan
[[Totale Differenzierbarkeit/R/Produktregel/Fakt|Lemma 45.10 (Analysis (Osnabrück 2021-2023))]]
berlaku

\mavergleichskettealignhandlinks
{\vergleichskettealignhandlinks
{ \left\langle  v  ,  L_P(w)  \right\rangle
}
{ =} { -\left\langle  v  ,  \left(  D_{w}  N   \right) \left(   P   \right)   \right\rangle
}
{ =} { -\left\langle  v  ,  \left(  D_{w}  { \frac{
\operatorname{Grad} \,  h  }{  \Vert {
\operatorname{Grad} \,  h } \Vert  } }   \right) \left(   P   \right)     \right\rangle
}
{ =} { -\left\langle  v  ,  \left(  D_{w}  { \frac{   1  }{  \Vert {
\operatorname{Grad} \,  h } \Vert  } }   \right) \left(   P   \right) \cdot
\operatorname{Grad} \,  h  (  P ) + { \frac{   1  }{  \Vert {
\operatorname{Grad} \,  h  (  P ) } \Vert  } } \cdot  \left(  D_{w}
\operatorname{Grad} \,  h   \right) \left(   P   \right)   \right\rangle
}
{ =} { - { \frac{   1  }{  \Vert {
\operatorname{Grad} \,  h  (  P ) } \Vert  } } \left\langle  v  ,  \left(  D_{w}
\operatorname{Grad} \,  h   \right) \left(   P   \right)    \right\rangle
}
}
{}
{}{,}
karena suku pertama tegak lurus terhadap vektor singgung $v$. Dalam fungsi-fungsi koordinat, berlaku

\mavergleichskettealignhandlinks
{\vergleichskettealignhandlinks
{ \left(  D_{w}
\operatorname{Grad} \,  h   \right) \left(   P   \right)
}
{ =} { \left(  D_{w}  \left(  { \frac{   \partial   h   }{  \partial   x_1   } }   , \,   \ldots  , \,   { \frac{   \partial   h   }{  \partial   x_n   } }                 \right)    \right) \left(   P   \right)
}
{ =} { \sum_{j = 1}^n  w_j  { \frac{   \partial    }{  \partial   x_j   } }  \left(  { \frac{   \partial   h   }{  \partial   x_1   } }   , \,   \ldots  , \,   { \frac{   \partial   h   }{  \partial   x_n   } }                 \right) (P)
}
{ =} { \left(  \sum_{j = 1}^n  w_j  { \frac{   \partial    }{  \partial   x_j   } }  { \frac{   \partial   h   }{  \partial   x_1   } } (P)   , \,   \ldots  , \,   \sum_{j = 1}^n  w_j  { \frac{   \partial    }{  \partial   x_j   } }  { \frac{   \partial   h   }{  \partial   x_n   } } (P)                 \right)
}
{ } {
}
}
{}
{}{.}
Dengan demikian, ekspresi di atas sama dengan

\mavergleichskettealign
{\vergleichskettealign
{ \left\langle  v  ,  L_P(w)  \right\rangle
}
{ =} { - { \frac{   1  }{  \Vert {
\operatorname{Grad} \,  h  (  P ) } \Vert  } } \left\langle  v  ,  \left(  D_{w}
\operatorname{Grad} \,  h   \right) \left(   P   \right)   \right\rangle
}
{ =} { -  { \frac{   1  }{  \Vert {
\operatorname{Grad} \,  h  (  P ) } \Vert  } }  \sum_{i,j = 1}^n v_i  w_j  { \frac{   \partial    }{  \partial   x_j   } } { \frac{   \partial   h   }{  \partial   x_i   } }  (P)
}
{ } {
}
{ } {
}
}
{}
{}{.}
Berdasarkan
[[Differenzierbarkeit/Satz von Schwarz/Fakt|Satz 44.10 (Analysis (Osnabrück 2021-2023))]]
urutan turunan parsial dapat dipertukarkan, sehingga peran
\mathkor {} {v} {dan} {w} {}
juga dapat dipertukarkan.

 }

__NOEDITSECTION__

\inputaufgabeklausurloesung
{4}

{

Misalkan $X$ adalah suatu
\definitionsverweis {ruang kompak}{}{}
dan misalkan

\mavergleichskette
{\vergleichskette
{Y
}
{ \subseteq }{ X
}
{  }{
}
{    }{
}
{   }{
}
}
{}{}{}
adalah
\definitionsverweis {himpunan bagian tertutup}{}{,}
yang dilengkapi dengan
\definitionsverweis {topologi terinduksi}{}{.}
Tunjukkan bahwa $Y$ juga kompak.

}
{

Misalkan

\mavergleichskettedisp
{\vergleichskette
{ Y
}
{ =} { \bigcup_{i \in I} V_i
}
{ } {
}
{ } {
}
{ } {
}
}
{}{}{}
adalah suatu tutupan oleh himpunan-himpunan bagian terbuka

\mavergleichskette
{\vergleichskette
{ V_i
}
{ \subseteq }{ Y
}
{  }{
}
{    }{
}
{   }{
}
}
{}{}{.}
Ini berarti bahwa terdapat himpunan-himpunan terbuka

\mavergleichskette
{\vergleichskette
{ U_i
}
{ \subseteq }{ X
}
{  }{
}
{    }{
}
{   }{
}
}
{}{}{}
yang memenuhi

\mavergleichskette
{\vergleichskette
{ V_i
}
{ = }{ Y \cap U_i
}
{  }{
}
{    }{
}
{   }{
}
}
{}{}{.}
Oleh karena itu,

\mavergleichskettedisp
{\vergleichskette
{ Y
}
{ \subseteq} { \bigcup_{i \in I} U_i
}
{ } {
}
{ } {
}
{ } {
}
}
{}{}{.}
Karena $Y$ tertutup dalam $X$, himpunan $X \setminus Y$ terbuka, sehingga

\mavergleichskettedisp
{\vergleichskette
{ X
}
{ =} { Y \cup (X \setminus Y)
}
{ =} { \bigcup_{i \in I} U_i  \cup  (X \setminus Y)
}
{ } {
}
{ } {
}
}
{}{}{}
adalah suatu tutupan terbuka dari $X$. Karena $X$ kompak, terdapat suatu subtutupan berhingga

\mavergleichskettedisp
{\vergleichskette
{ X
}
{ =} { \bigcup_{i \in J} U_i  \cup  (X \setminus Y)
}
{ } {
}
{ } {
}
{ } {}
}
{}{}{.}
Ini selanjutnya berarti

\mavergleichskettedisp
{\vergleichskette
{ Y
}
{ =} { \bigcup_{i \in J} U_i
}
{ } {
}
{ } {
}
{ } {}
}
{}{}{.}

 }

__NOEDITSECTION__

\inputaufgabeklausurloesung
{5}

{

Berikan contoh suatu
\definitionsverweis {manifold diferensiabel}{}{}
$M$ yang
\definitionsverweis {berdimensi dua}{}{}
dan
\definitionsverweis {terhubung}{}{,}
beserta suatu titik

\mavergleichskette
{\vergleichskette
{ P
}
{ \in }{ M
}
{  }{
}
{    }{
}
{   }{
}
}
{}{}{}
sedemikian sehingga
\mathkor {} {M} {dan} {M \setminus \{P\}} {}
saling
\definitionsverweis {difeomorfik}{}{.}

}
{

Misalkan

\mavergleichskettedisp
{\vergleichskette
{M
}
{ =} { \R^2 \setminus \N_+
}
{ =} { \R^2 \setminus \{ (1,0) , (2,0), (3,0), \ldots  \}
}
{ } {
}
{ } {
}
}
{}{}{,}
yakni dari $\R^2$ dihapus bilangan-bilangan asli positif yang terletak pada sumbu $x$. Himpunan ini merupakan himpunan bagian terbuka dari $\R^2$, sehingga merupakan manifold diferensiabel. Manifold ini terhubung lintasan karena, misalnya, setiap titik di $M$ dengan
\mathl{y \geq 0}{} dapat dihubungkan dengan
\mathl{(0,1)}{} melalui lintasan garis lurus dan setiap titik di $M$ dengan
\mathl{y \leq 0}{} dapat dihubungkan dengan
\mathl{(0,-1)}{} melalui lintasan garis lurus, sedangkan kedua titik tersebut juga dapat dihubungkan dengan garis lurus. Sekarang misalkan

\mavergleichskettedisp
{\vergleichskette
{P
}
{ =} { (0,0)
}
{ } {
}
{ } {
}
{ } {
}
}
{}{}{}
dan

\mavergleichskettedisp
{\vergleichskette
{M'
}
{ =} { M \setminus \{P\}
}
{ } {
}
{ } {
}
{ } {
}
}
{}{}{.}
Pemetaan
\maabbeledisp {} {\R^2} {\R^2
} {(x,y)} {(x-1,y)
} {,}
merupakan suatu difeomorfisme
\zusatzklammer {sebagai translasi} {} {}
dan citra $M$ tepat sama dengan $M'$. Oleh karena itu,
\mathkor {} {M} {dan} {M'=M \setminus \{P\}} {}
saling difeomorfik.

 }

__NOEDITSECTION__

\inputaufgabeklausurloesung
{3}

{

Untuk setiap
\definitionsverweis {vektor singgung}{}{}

\mavergleichskette
{\vergleichskette
{ v
}
{ \in }{ T_PS^2
}
{ \subseteq }{ \R^3
}
{    }{
}
{   }{
}
}
{}{}{}
dengan

\mavergleichskette
{\vergleichskette
{ \Vert {v} \Vert
}
{ = }{ 1
}
{  }{
}
{    }{
}
{   }{
}
}
{}{}{}
di suatu titik

\mavergleichskette
{\vergleichskette
{ P
}
{ \in }{ S^2
}
{ \subset }{ \R^3
}
{    }{
}
{   }{
}
}
{}{}{}
pada
\definitionsverweis {permukaan bola satuan}{}{,}
berikan suatu wakil
\definitionsverweis {diferensiabel}{}{}
\maabbdisp {\gamma} { I } { S^2
} {}
dengan

\mavergleichskette
{\vergleichskette
{ \gamma(0)
}
{ = }{ P
}
{  }{
}
{    }{
}
{   }{
}
}
{}{}{.}

}
{

Misalkan

\mavergleichskette
{\vergleichskette
{P
}
{ = }{ \begin{pmatrix}  x  \\ y \\ z   \end{pmatrix}
}
{  }{
}
{    }{
}
{   }{
}
}
{}{}{}
dan

\mavergleichskette
{\vergleichskette
{v
}
{ = }{ \begin{pmatrix}  a  \\ b \\ c   \end{pmatrix}
}
{  }{
}
{    }{
}
{   }{
}
}
{}{}{.}
Kedua vektor tersebut bernorma satu dan saling tegak lurus karena bidang singgung pada permukaan bola tegak lurus terhadap vektor posisi. Tinjau lintasan
\maabbdisp {\gamma} {\R} { \R^3
} {}
yang diberikan oleh

\mavergleichskettedisp
{\vergleichskette
{ \gamma(t)
}
{ =} {  \left(  \cos t  \right)  \begin{pmatrix}  x  \\ y \\ z   \end{pmatrix}  + \left(  \sin t  \right)  \begin{pmatrix}  a  \\ b \\ c   \end{pmatrix}
}
{ } {
}
{ } {
}
{ } {
}
}
{}{}{.}
Berlaku

\mavergleichskettedisp
{\vergleichskette
{ \Vert { \gamma(t)} \Vert
}
{ =} { \cos^{ 2  }  t + \sin^{ 2  }  t
}
{ =} { 1
}
{ } {
}
{ } {
}
}
{}{}{,}
sehingga lintasan tersebut seluruhnya berada pada permukaan bola. Selanjutnya,

\mavergleichskettedisp
{\vergleichskette
{ \gamma(0)
}
{ =} { \left(  \cos  0  \right) \begin{pmatrix}  x  \\ y \\ z   \end{pmatrix} + \left(  \sin  0  \right)  \begin{pmatrix}  a  \\ b \\ c   \end{pmatrix}
}
{ =} { \begin{pmatrix}  x  \\ y \\ z   \end{pmatrix}
}
{ } {
}
{ } {
}
}
{}{}{}
dan

\mavergleichskettedisp
{\vergleichskette
{ \gamma'(0)
}
{ =} { - \left(  \sin  0  \right) \begin{pmatrix}  x  \\ y \\ z   \end{pmatrix} + \left(  \cos  0  \right) \begin{pmatrix}  a  \\ b \\ c   \end{pmatrix}
}
{ =} { \begin{pmatrix}  a  \\ b \\ c   \end{pmatrix}
}
{ } {
}
{ } {
}
}
{}{}{.}

 }

__NOEDITSECTION__

\inputaufgabeklausurloesung
{2}

{

Tinjau parabola standar

\mavergleichskette
{\vergleichskette
{ Y
}
{ \subseteq }{ \R^2
}
{  }{
}
{    }{
}
{   }{
}
}
{}{}{}
dengan struktur Riemann terinduksi dan parametrisasi
\maabbeledisp {\varphi} {\R} { Y
} { t } { \left(  t   , \,  t^2                   \right)
} {,}
yang kita pandang sebagai suatu peta
\zusatzklammer {invers} {} {}.
Tentukan
\definitionsverweis {fungsi fundamental}{}{}
Riemann

\mavergleichskette
{\vergleichskette
{ g(t)
}
{ = }{ g_{11} (t)
}
{  }{
}
{    }{
}
{   }{
}
}
{}{}{.}

}
{

Berlaku

\mavergleichskettealign
{\vergleichskettealign
{ g(t)
}
{ =} { \left\langle  T_t \varphi(1)   , T_t \varphi(1)   \right\rangle_{ Y, \varphi(t) }
}
{ =} { \left\langle    \varphi' (t)   ,   \varphi'(t)   \right\rangle
}
{ =} { \left\langle  \begin{pmatrix}  1 \\ 2t   \end{pmatrix}     ,  \begin{pmatrix}  1 \\ 2t   \end{pmatrix}   \right\rangle
}
{ =} { 1+4t^2
}
}
{}
{}{.}

 }

__NOEDITSECTION__

\inputaufgabeklausurloesung
{6}

{

Hitung bentuk diferensial tarik balik
\mathl{\varphi^* \tau}{} dari

\mavergleichskettedisp
{\vergleichskette
{ \tau
}
{ =} { dx \wedge dy \wedge dz - w dx \wedge dy \wedge dw + \cos (xy)  dx \wedge dz \wedge dw-yw dy \wedge dz \wedge dw
}
{ } {
}
{ } {
}
{ } {
}
}
{}{}{}
oleh pemetaan
\maabbeledisp {\varphi} { \R^3 } { \R^4
} { (r,s,t) } {  \left(  r^2s,t, \sin  r ,e^{st}  \right)  = (x,y,z,w)
} {.}

}
{

Diperoleh

\mavergleichskettedisp
{\vergleichskette
{dx
}
{ =} { d(r^2s)
}
{ =} { 2rs dr  +r^2ds
}
{ } {
}
{ } {
}
}
{}{}{,}

\mavergleichskettedisp
{\vergleichskette
{dy
}
{ =} { dt
}
{ } {
}
{ } {
}
{ } {
}
}
{}{}{,}

\mavergleichskettedisp
{\vergleichskette
{dz
}
{ =} { d (\sin  r )
}
{ =} { \cos  r  dr
}
{ } {
}
{ } {
}
}
{}{}{}
dan

\mavergleichskettedisp
{\vergleichskette
{dw
}
{ =} { d(e^{st})
}
{ =} { se^{st} dt+te^{st} ds
}
{ } {
}
{ } {
}
}
{}{}{.}
Dengan demikian,

\mavergleichskettealign
{\vergleichskettealign
{ \varphi^* \tau
}
{ =} { d(r^2s) \wedge dt \wedge d( \sin  r )  -e^{st} d(r^2s) \wedge dt \wedge d(e^{st})
}
{ \, \, \, \,} {+ \cos (r^2st)  d (r^2s) \wedge d ( \sin  r ) \wedge d(e^{st}) - te^{st} dt \wedge d( \sin  r ) \wedge d(e^{st})
}
{ =} { r^2 \cos (r^2st) \cos  r  ds \wedge dt \wedge dr  - 2rs t e^{st} e^{st} dr \wedge dt \wedge ds
}
{ \, \, \, \,} { + r^2 se^{st} \cos  r  ds \wedge dr \wedge dt  -  t^2 e^{st}  e^{st} \cos  r dt \wedge dr \wedge ds
}
}
{
\vergleichskettefortsetzungalign
{ =} { (r^2 \cos  r +2 rst e^{2st} - r^2 s e ^{st} \cos (r^2st) \cos  r - t^2e^{ 2st} \cos  r     )  dr \wedge ds \wedge dt
}
{ } {}
{ } {}
{ } {}
}
{}{}

 }

__NOEDITSECTION__

\inputaufgabeklausurloesung
{7 (2+3+2)}

{

Tinjau pemetaan-pemetaan diferensiabel
\maabbeledisp {\gamma} { [1,2] } {\R^2
} { t } {(t,t^{-1})
} {,}
dan
\maabbeledisp {\varphi} {\R^2} {\R^3
} {(u,v)} {(u^2,uv,-u+v^2)
} {,}
serta bentuk diferensial

\mavergleichskettedisp
{\vergleichskette
{ \omega
}
{ =} { xdx-zdy+dz
}
{ } {
}
{ } {
}
{ } {
}
}
{}{}{}
pada $\R^3$.

a) Hitung bentuk diferensial tarik balik
\mathl{\varphi^* \omega}{} pada $\R^2$.

b) Hitung integral lintasan bentuk diferensial
\mathl{\varphi^* \omega}{} sepanjang lintasan $\gamma$.

c) Hitung
\zusatzklammer {tanpa menggunakan hasil b)} {} {}
integral lintasan bentuk diferensial
\mathl{\omega}{} sepanjang lintasan $\varphi \circ \gamma$.

}
{

a) Bentuk diferensial tarik baliknya adalah

\mavergleichskettealign
{\vergleichskettealign
{\varphi^* ( xdx -zdy +dz )
}
{ =} { u^2 du^2    -(-u+v^2) duv +d (-u+v^2)
}
{ =} { 2u^3 du + (u-v^2) (vdu +udv) -du +2vdv
}
{ =} { (2u^3 +uv -v^3 -1) du    + (u^2 -uv^2 +2v )dv
}
{ } {
}
}
{}
{}{.}

b) Integral lintasannya adalah

\mavergleichskettealign
{\vergleichskettealign
{ \int_1^2 ( 2t^3 +1 -t^{-3} -1 )dt  +(t^2 -t^{-1} +2t^{-1} ) d t^{-1}
}
{ =} { \int_1^2 ( 2t^3 -t^{-3} ) dt  +(t^2 + t^{-1} ) d t^{-1}
}
{ =} { \int_1^2 ( 2t^3 -t^{-3} ) dt  -(1 + t^{-3} ) dt
}
{ =} { \int_1^2 ( 2t^3 - 2t^{-3} -1 ) dt
}
{ =} { \left(  { \frac{   1 }{  2 } }  t^4  + t^{-2}  -t  \right) \vert_{  1 }^{  2 }
}
}
{
\vergleichskettefortsetzungalign
{ =} { 8 + { \frac{   1 }{  4 } }  -2 - { \frac{   1 }{  2 } } -1 +1
}
{ =} { 5 + { \frac{   3 }{  4 } }
}
{ } {}
{ } {}
}
{}{.}

c) Lintasan komposisinya adalah

\mavergleichskettedisp
{\vergleichskette
{ \psi (t)
}
{ =} { (t^2,1,-t +t^{-2} )
}
{ } {
}
{ } {
}
{ } {
}
}
{}{}{.}
Dengan demikian,

\mavergleichskettealign
{\vergleichskettealign
{ \int_1^2 \psi^* (xdx -zdy +dz )
}
{ =} { \int_1^2 t^2dt^2 + d (-t + t^{-2} )
}
{ =} {  \int_1^2 (2 t^3-1 -2t^{-3} ) dt
}
{ =} { 5 + { \frac{   3 }{  4 } }
}
{ } {
}
}
{}
{}{.}

 }

__NOEDITSECTION__

\inputaufgabeklausurloesung
{4}

{

Misalkan $M$ adalah suatu
\definitionsverweis {manifold berbatas}{}{.}
Tunjukkan bahwa setiap himpunan bagian terbuka

\mavergleichskette
{\vergleichskette
{ N
}
{ \subseteq }{ M
}
{  }{
}
{    }{
}
{   }{
}
}
{}{}{}
juga merupakan manifold dengan batas
\zusatzklammer {yang mungkin kosong} {} {},
dan bahwa

\mavergleichskettedisp
{\vergleichskette
{ \partial N
}
{ =} { \partial M \cap N
}
{ } {
}
{ } {
}
{ } {
}
}
{}{}{}
berlaku.

}
{

Misalkan
\mathl{P \in N}{} dan
\mathl{P \in U}{} adalah suatu domain peta dari $M$. Misalkan
\maabbdisp {\alpha} { U } {V
} {}
adalah suatu peta dengan himpunan terbuka
\mathl{V \subseteq H}{,}
\mathl{H=\R_{\geq 0} \times \R^{n-1}}{.} Peta ini menginduksi suatu homeomorfisme
\maabbdisp {} {N \cap U} { \alpha (N \cap U)
} {.}
Peta-peta ini kita gunakan untuk $N$. Sifat difeomorfisme dari transisi peta pada $M$ langsung diwarisi oleh peta-peta pada $N$. Misalkan
\mathl{P \in N \cap \partial M}{.} Maka suatu peta memetakan $P$ ke titik batas $H$. Karena peta-peta untuk $N$ merupakan pembatasan dari peta-peta tersebut, sifat ini tetap berlaku. Sebaliknya, jika
\mathl{P \in \partial N}{,} maka di satu pihak tentu saja
\mathl{P \in N \subseteq M}{.} Di pihak lain, untuk $P$ terdapat suatu peta pada $N$ yang diperoleh dengan membatasi suatu peta untuk $M$, dan peta tersebut memetakan $P$ ke titik batas $H$.

 }

__NOEDITSECTION__

\inputaufgabeklausurloesung
{8}

{

Buktikan teorema tentang karakterisasi kurva geodesik terhadap koneksi Levi-Civita.

}
{

Misalkan $n$ adalah dimensi $M$. Tinjau situasi ini secara langsung pada suatu citra peta terbuka

\mavergleichskette
{\vergleichskette
{ U
}
{ \subseteq }{ \R^n
}
{  }{
}
{    }{
}
{   }{
}
}
{}{}{.}
Menurut
[[Offene Menge/Triviales Vektorbündel/Linearer Zusammenhang/Spaltung/Christoffelsymbole/Bemerkung|Bemerkung]] [[Kurs:Differentialgeometrie/Offene Menge/Triviales Vektorbündel/Linearer Zusammenhang/Spaltung/Christoffelsymbole/Bemerkung/Bemerkungreferenznummer|*****]]
turunan vertikal diberikan oleh
\maabbeledisp {} {T(U \times \R^n) = U \times \R^n \times \R^n \times \R^n } { U \times \R^n \times  \R^n
} {(P,e_j, \partial_i, w)} {(P,e_j,\sum_{k = 1}^n \Gamma^k_{ij} (P) e_k +w) \cong V(U \times \R^n)
} {}
; pemetaan ke $TU$ diperoleh dengan menghilangkan komponen tengah. Turunan kedua kurva mula-mula merupakan pemetaan singgung kedua, yaitu
\zusatzklammer {di sini kita mengabaikan perkalian dengan arah satu dimensi dari ruang singgung kurva} {} {}
\maabbeledisp {T(\gamma)} { I } { TU
} { t } { (\gamma(t), \gamma'(t))
} {,}
dan dengan demikian
\maabbeledisp {T(T(\gamma))} { I } { T(TU)
} { t } { ((\gamma(t), \gamma'(t)), (\gamma'(t), \gamma^{\prime \prime} (t)))
} {}
\zusatzklammer {kedua komponen pemetaan singgung diturunkan} {} {.}
Dengan

\mavergleichskette
{\vergleichskette
{ \gamma'(t)
}
{ = }{ \sum_{j = 1}^n \gamma_j'(t) \partial_j
}
{  }{
}
{    }{
}
{   }{
}
}
{}{}{}
proyeksi vertikal memetakan bentuk ini ke

\mathdisp {\sum_{k = 1}^n  \left(  \sum_{i,j} \gamma_i'(t) \gamma_j'(t) \Gamma^k_{ij}(\gamma(t))  \right)  \partial_k  +\sum_{k = 1}^n \gamma_k^{\prime \prime}(t) \partial_k} {  }

. Syarat bagi suatu kurva geodesik bahwa turunan keduanya
\zusatzklammer {di $TTM$} {} {}
selalu horizontal ekuivalen dengan proyeksi vertikal yang dihitung bernilai $0$. Ini berarti bahwa setiap komponennya bernilai $0$, yaitu

\mavergleichskettedisp
{\vergleichskette
{ \sum_{i,j} \gamma_i'(t) \gamma_j'(t) \Gamma^k_{ij}(\gamma(t))  + \gamma_k^{\prime \prime}(t)
}
{ =} { 0
}
{ } {
}
{ } {
}
{ } {
}
}
{}{}{}
untuk

\mavergleichskette
{\vergleichskette
{ k
}
{ = }{ 1  , \ldots , n
}
{  }{
}
{    }{
}
{   }{
}
}
{}{}{.}

 }

 [[Kategorie:Latexseite]]
